\section{相关工作}

\subsection{大语言模型辅助硬件形式化验证}

现有大语言模型辅助硬件形式化验证的研究，首先可以按其生成或处理的对象区分。\emph{AssertLLM} 从设计规格文档提取结构信息并生成待验证的硬件断言\cite{fang2024assertllm}；\emph{FVEval} 则以自然语言到 SVA、设计到 SVA 等任务评估模型生成断言的能力\cite{kang2024fveval}。这类工作主要解决\emph{验证目标如何表达}的问题。另有 \emph{FVDebug} 利用失败轨迹、RTL 和规格定位形式化验证失败的原因，并辅助修改设计\cite{bai2025fvdebug}；其任务是解释或修复失败的实现，而非为既定目标选择证明引理。

与本文更接近的是 \emph{Large Lemma Miners}：它让 LLM 提出 RTL 归纳证明的候选引理，再由形式化后端确认能否构成证明\cite{peled2026large}。该工作同时包含非 Agentic 与 Agentic 候选生成流程；后者会根据候选是否成立等检查结果继续生成引理。因此，\emph{形式化反馈指导 LLM} 已有先例。本文沿用固定的候选生成设置，研究候选池形成之后的组合选择，避免将生成策略的差异计入选择算法的效果。

\subsection{归纳引理生成与不变式推断}

利用验证反馈构造不变式也有较长的研究脉络。\emph{ICE} 用正例、反例和蕴含关系约束不变式学习\cite{garg2014ice}，表明反例不仅能否定候选，还能提供后续搜索所需的结构化约束。在软件验证中，Kamath 等人让 LLM 生成循环不变式，并用符号工具检查\cite{kamath2023loop}；\emph{Lemur} 让 LLM 提议证明子目标，再由自动推理器验证推理步骤\cite{wu2024lemur}。这些工作说明模型提议与形式化判定可以分工，但其主要对象是程序性质或循环不变式的生成，而非给定 RTL 候选池的低成本选择。

在候选组合方面，\emph{Clause2Inv} 让 LLM 生成程序循环不变式的子句，由符号组合器利用历史反例构造完整不变式；违反的证明义务还会引导后续子句生成\cite{cao2025clause2inv}。在反馈修复方面，\emph{LORIS} 检查 LLM 为失败证明义务给出的逐步推理，定位其中的局部逻辑错误，再引导模型重生成循环不变式\cite{li2026loris}。因此，反例驱动组合与结构化反馈并非本文独有的思路。

与上述工作相比，本文的决策对象是生成后冻结的 RTL 引理池。LLM 仅推荐池内引理 ID 的联合组合；形式化后端若返回可可靠求值的初始性或归纳步骤反例，选择器据此过滤或修复\emph{下一组待查组合}。该反馈不产生新引理，也不核查模型的自然语言推理链。模型请求、形式化查询及回退共用墙钟预算。

\subsection{前缀搜索与 Houdini 引理筛选}

在本文给定候选引理池的设置下，\emph{Houdini} 从全部候选出发，依据验证结果迭代剔除不满足初始性或归纳保持条件的候选，直到剩余集合不再变化\cite{flanagan2001houdini}。在候选语义一致且相关查询完成时，它可得到固定池内的最大有效归纳子集。另一种策略是先确定检查顺序：\emph{Large Lemma Miners} 将候选按归纳性与有界检查结果分类排序，再依次验证该顺序产生的前缀\cite{peled2026large}。两种方法采用不同的候选筛选策略。本文在相同候选池、验证条件和预算下，将其作为基线比较证明完成率与检查成本。

后续研究已经把目标信息和搜索成本纳入不变式求解。\emph{H-Houdini} 将硬件不变式学习中的整体归纳检查分解为性质导向、可增量和可复用的层次化检查\cite{dinesh2025hhoudini}；\emph{PyCaliper} 的 RTL 不变式合成框架允许指定候选搜索策略，并在缓存隔离的案例中让 LLM 根据当前不变式和目标性质排列候选顺序\cite{godbole2025pycaliper}。因此，性质导向筛选、增量验证和 LLM 候选排序都已有研究基础。本文关注的是在相同形式化后端和固定候选池上，由 LLM 提议候选 ID 组合，并在可靠的归纳反馈下调整后续待查组合，同时把模型调用与形式化查询纳入同一预算；H-Houdini 的检查分解与并行优化作为相关研究讨论，不作为本文的第三个主外部基线。

\subsection{本文方法定位}

综上，相关工作分别推进了断言或引理生成、反例指导的不变式合成，以及固定候选的排序与筛选。本文选择其中一个更受限的决策点：对同一 RTL 设计、目标属性及已冻结的候选池，预测可能相互支持的非前缀组合；若形式化检查失败，仅在反例语义可靠时据此修复或排除组合，并以统一时间预算决定继续搜索还是回退。形式化后端始终负责初始性与归纳性的最终判定。

本文以相同候选池下的前缀搜索和经典 Houdini 为基线，检验组合选择的收益能否覆盖 LLM 开销；同一次模型输出的排序前缀、以及仅接收成败信息的重试，分别用于检验组合级选择与结构化反馈的增量作用。
